% 2026-09-11：非数学模式的西文字母统一使用 Pagella 正体，保留粗细区别。
% 保留字母级边界，不重分类数字、括号、标点及中文。
% 此设置也覆盖列表序号字母和术语原文；数学模式保持原样。
\newfontfamily\AnalysisLetterFont{texgyrepagella-regular.otf}[
  BoldFont=texgyrepagella-bold.otf,
  ItalicFont=texgyrepagella-italic.otf,
  BoldItalicFont=texgyrepagella-bolditalic.otf
]
\ExplSyntaxOn
\bool_new:N \l_analysis_text_letters_bool
\bool_set_true:N \l_analysis_text_letters_bool
\cs_new_protected:Npn \AnalysisPreserveMathText
  { \bool_set_false:N \l_analysis_text_letters_bool }
\cs_new_protected:Npn \analysis_letter_group_begin:
  {
    \group_begin:
    \bool_if:NT \l_analysis_text_letters_bool
      { \AnalysisLetterFont \upshape }
  }
% 继承 Default 的全部 xeCJK 转移，保留中西文间距及标点处理；
% 不能用只有开关分组的转移覆盖 Boundary 或 CJK 的原处理。
\xeCJK_new_class:n { AnalysisLatinLetter }
\seq_map_inline:Nn \g__xeCJK_class_seq
  {
    \str_if_eq:nnF {#1} { AnalysisLatinLetter }
      {
        \xeCJK_copy_inter_class_toks:nnnn
          {#1} {AnalysisLatinLetter} {#1} {Default}
        \xeCJK_copy_inter_class_toks:nnnn
          {AnalysisLatinLetter} {#1} {Default} {#1}
        \xeCJK_app_inter_class_toks:nnn {#1} {AnalysisLatinLetter}
          { \analysis_letter_group_begin: }
        \xeCJK_pre_inter_class_toks:nnn {AnalysisLatinLetter} {#1}
          { \group_end: }
      }
  }
% ASCII 字母及常用拉丁扩展字母；明确跳过乘号、除号，不重分类标点。
\xeCJKDeclareCharClass { AnalysisLatinLetter }
  { "0041->"005A, "0061->"007A, "00C0->"00D6, "00D8->"00F6,
    "00F8->"024F, "1E00->"1EFF }
\ExplSyntaxOff
% 数学组内局部关闭字母切换，连公式中 \text{...} 的字体也保持原样。
% 只追加，不覆盖 amsmath、unicode-math、xeCJK 已有的数学初始化。
\AtBeginDocument{%
  \everymath\expandafter{\the\everymath\AnalysisPreserveMathText}%
  \everydisplay\expandafter{\the\everydisplay\AnalysisPreserveMathText}%
}
